\section{Introduction: designing a system that delivers the water}
The first five chapters answered how much water a soil holds, how it moves, when to irrigate and how much a crop uses. A farmer has still to get that water to the plant: from a well or a canal, through pipes or channels, into the soil, with what pressure and at what cost. This is \kw{irrigation design}. It rests on a few ideas from hydraulics (continuity, energy, friction), on a clear idea of efficiency (how much of the water taken from the source reaches the roots), and on data about the weather and the field that set the demand. This chapter collects these ideas; the later chapters apply them to surface, drip and sprinkler systems and to drainage.

\section{Classification of irrigation methods}
Irrigation methods divide into three families by the way the water reaches the soil. In \kw{surface irrigation} the water flows over the field by gravity: \emph{basin} irrigation floods a levelled area, \emph{border} irrigation uses long strips with a free downstream end, and \emph{furrow} irrigation runs the water in small channels between the crop rows. It needs no pressure, but the application is non-uniform unless the field is well levelled and the inflow well chosen; efficiencies are 50--70 per cent. In \kw{sprinkler irrigation} water is sprayed through nozzles under pressure: \emph{set} systems (solid set, hand-move, side-roll), \emph{moving} systems (centre pivot, lateral move, travelling gun) and \emph{rotating} or \emph{impact} sprinklers differ in cost and labour; efficiencies are 70--85 per cent. In \kw{micro-irrigation} (drip, micro-sprinkler, subsurface drip) small flows are given to each plant through emitters at low pressure; efficiencies are 85--95 per cent and the wetted volume is a fraction of the soil, which suits row crops, trees and vines and needs filtration to prevent clogging.

\section{Efficiency and uniformity}
Not all water taken from the source benefits the crop. Four efficiencies are used. The \kw{conveyance efficiency} $E_c$ is the ratio of the water that reaches the field to the water taken at the source; the \kw{application efficiency} $E_a=V_{stored}/V_{applied}$ is the share of the applied water stored in the root zone; the \kw{storage efficiency} $E_s=V_{stored}/V_{needed}$ is the share of the needed depth that has been supplied; and the \kw{distribution uniformity} $DU_{lq}$ is the mean depth of the lowest quarter divided by the mean depth. The overall field efficiency is $E=E_cE_a$. The \kw{Christiansen coefficient} $CU=1-\sum|z_i-\bar z|/(n\bar z)$ is the other standard measure of uniformity, and, for emitters, the \kw{emission uniformity} $EU=q_{min}/\bar q$.

\exbox{Uniformity from catch cans}{Twelve catch cans in a sprinkler field collect (mm): 12, 14, 15, 15, 16, 17, 17, 18, 19, 22, 10, 13. Find the mean, Christiansen's coefficient, and the distribution uniformity of the lowest quarter.}\sol Sum $=188$, mean $\bar z=15.67$ mm. Absolute deviations sum to $30.00$; $CU=1-30.00/(12\times15.67)=0.840$ \ans Lowest quarter: the three smallest values $10, 12, 13$, mean $11.67$; $DU_{lq}=11.67/15.67=0.745$ \finans

\section{Hydraulics: continuity, energy and friction}
Pressurised systems obey three relations. \kw{Continuity}: the flow $Q=vA$ is the same at every section of a pipe without branches, so a pipe of diameter $D$ carries the velocity $v=4Q/(\pi D^2)$. \kw{Energy} (Bernoulli): the total head $H=z+p/(\rho g)+v^2/(2g)$ falls along a pipe by the friction head loss $h_f$ and the local losses $h_m=K_mv^2/(2g)$ of fittings. \kw{Friction}: for water in a smooth plastic pipe the empirical Hazen--Williams equation in SI units is

\begin{equation}h_f=10.67\,\frac{L\,Q^{1.852}}{C^{1.852}\,D^{4.871}},\label{eq:hw6}\end{equation}


with $Q$ in m$^3$\,s$^{-1}$, $D$ in m, $L$ in m and $C$ the roughness coefficient (140 for new PVC and polyethylene pipe, 150 for the smooth laterals; lower for old steel). The Darcy--Weisbach form $h_f=f\,(L/D)\,v^2/(2g)$ with the friction factor $f$ from the Moody diagram is more general. A pipe is sized so that the velocity lies between 0.3 m\,s$^{-1}$ (to keep sediment moving) and 1.5 m\,s$^{-1}$ (to limit water hammer), and so that the head loss in the pipe and the cost of the pipe together are the least. The power of a pump that delivers $Q$ at a total dynamic head $H$ with an efficiency $\eta$ is $P=\rho g Q H/\eta$.

\exbox{Mainline friction and power}{A PVC mainline of $L=600$ m and $D=200$ mm ($C=140$) carries $Q=45$ L\,s$^{-1}$. Find the velocity, the head loss, and the extra pump power needed to overcome it ($\eta=0.65$).}\sol \givenline{$L=600$ m ; $D=200$ mm ; $C=140$ ; $Q=45$ L\,s$^{-1}$ ; $\eta=0.65$}\ans $Q=0.045$ m$^3$\,s$^{-1}$; $v=4\times0.045/(\pi\times0.2^2)=1.43$ m\,s$^{-1}$ (within 0.3--1.5). Head loss: $10.67\times600\times0.045^{1.852}/(140^{1.852}\times0.2^{4.871})=5.52$ m \ans Power $=\rho gQh_f/\eta=1000\times9.81\times0.045\times5.52/0.65=3750$ W $=3.75$ kW \finans

\section{Weather data for a design}
Every design and every model needs the weather: rainfall, maximum and minimum temperature, solar radiation (or sunshine hours), humidity (or dew point) and wind speed, daily, for many years. The data may come from a weather station near the field, from the national service, from a gridded reanalysis or from a satellite product. Three checks are needed before use: \emph{completeness} (gaps are filled by interpolation from neighbouring days or stations), \emph{plausibility} (negative rain, $T_{min}>T_{max}$, radiation above the extraterrestrial value are errors) and \emph{homogeneity} (a change of instrument shows as a step in the mean). When radiation is missing it may be estimated from the temperature range by Hargreaves' relation $R_s=k_{Rs}\sqrt{T_{max}-T_{min}}\,R_a$ with $k_{Rs}\approx0.16$ (interior) to 0.19 (coastal) and $R_a$ the extraterrestrial radiation. For a design the \emph{dependable} (e.g. 80 per cent probability) rainfall and the \emph{peak} $ET_0$ of the design period are used, not the means.

\exbox{Radiation from the temperature range}{A station lacks radiation. For $T_{max}=31$, $T_{min}=17$ $^\circ$C and an extraterrestrial radiation $R_a=38$ MJ\,m$^{-2}$\,d$^{-1}$ (mid-latitude summer day), estimate $R_s$ with $k_{Rs}=0.17$ and the net shortwave radiation for an albedo 0.23.}\sol \givenline{$T_{max}=31$ ; $T_{min}=17$ $^\circ$C ; $R_a=38$ MJ\,m$^{-2}$\,d$^{-1}$ ; $k_{Rs}=0.17$}\ans $R_s=0.17\sqrt{31-17}\times38=0.17\times3.742\times38=24.2$ MJ\,m$^{-2}$\,d$^{-1}$ \ans $R_{ns}=(1-0.23)R_s=18.6$ MJ\,m$^{-2}$\,d$^{-1}$ \ans The clear-sky bound is $0.75R_a=28.5$, so the estimate is plausible. \finans

\section{Field data from satellites: vegetation indices}
A satellite measures the reflectance of the surface in bands of the spectrum. Healthy green vegetation absorbs red light (chlorophyll) and reflects near-infrared strongly (leaf structure); bare soil reflects both moderately; water absorbs both. The \kw{normalised difference vegetation index} is

\begin{equation}\mathrm{NDVI}=\frac{\rho_{NIR}-\rho_{red}}{\rho_{NIR}+\rho_{red}},\label{eq:ndvi6}\end{equation}
\par\noindent where NDVI is the normalised difference vegetation index, $\rho_{NIR}$ the reflectance in the near-infrared band and $\rho_{red}$ the reflectance in the red band (both 0 to 1).



from $-1$ (water) through 0 (bare soil) to about 0.9 (dense canopy). Other indices use other bands: the \emph{normalised difference water index} with a shortwave band follows the water in the canopy, and the \emph{soil-adjusted} index corrects for soil brightness. The fraction of the ground covered by the canopy is estimated from NDVI by $f_c=[(\mathrm{NDVI}-\mathrm{NDVI}_s)/(\mathrm{NDVI}_v-\mathrm{NDVI}_s)]^2$, and the basal crop coefficient from $f_c$; this is how an image can drive a schedule.

\exbox{NDVI and cover}{A pixel has reflectances $\rho_{red}=0.08$ and $\rho_{NIR}=0.42$. Find the NDVI and the cover fraction if $\mathrm{NDVI}_s=0.15$ (bare) and $\mathrm{NDVI}_v=0.85$ (full cover).}\sol \givenline{$\rho_{red}=0.08$ ; $\rho_{NIR}=0.42$ ; $\mathrm{NDVI}_s=0.15$ ; $\mathrm{NDVI}_v=0.85$}\ans $\mathrm{NDVI}=(0.42-0.08)/(0.42+0.08)=0.34/0.50=0.680$ \ans $f_c=[(0.680-0.15)/(0.85-0.15)]^2=(0.757)^2=0.573$ \ans About 57 per cent of the ground is covered. \finans

